\clearpage\SourceLeaf{401}\ArithmeticHead{394}{L}{ANALYSIS. NACHLASS.}
\noindent\hspace*{\parindent}Wenn beide Umstände sich vereinigen, dass die ursprüngliche Reihe sowohl eine obere als eine untere Grenze hat, oder zur IV\,ten Gattung gehört, so kann in [der] Reihe \(L',\,L'',\,L'''\) u. s. f. kein Glied vorkommen, das kleiner als irgend ein Glied der Reihe \(M',\,M'',\,M'''\) u. s. f. wäre. Denn gesetzt es wäre \(L'<M^m\), so würde wenn \(n\) irgend eine ganze Zahl bedeutet, die kleiner als 1 und \(m\) (oder auch der kleinsten von beiden gleich) ist, sein \(L^n\leqq L'\), \(M^n\geqq M^m\) und folglich \(L^n<M^n\), welches widersprechend ist, da \(L^n\) die obere, \(M^n\) die untere Grenze von Einer und derselben Reihe \(a^n,a^{n+1},a^{n+2}\) etc. bedeuten. — Hieraus folgt nun leicht, dass die Reihe \(L',\,L'',\,L'''\) u. s. f. eine untere, die Reihe \(M',\,M'',\,M'''\) u. s. f. aber eine obere Grenze haben werde; jene wollen wir durch \(L\), diese [durch] \(M\) bezeichnen\textsuperscript{*}), und man begreift leicht, dass \(L\) nicht kleiner als \(M\) sein könne, sondern entweder \(L>M\) oder \(L=M\) sein müsse. \(L\) nennen wir die letzte obere Grenze der ursprünglichen Reihe \([R]\), \(M\) die letzte untere derselben. Zugleich erhellet, dass \(L,\,M\) auch die letzten beiderseitigen Grenzen der Reihen \(a'',\,a''',\,a^{\mathrm{IV}}\) u. s. f.; \(a''',\,a^{\mathrm{IV}}\) u. s. f. u. s. f. sein werden, oder dass man bei Bestimmung der letzten Grenzen einer Reihe von ihren Anfangsgliedern so viele weglassen kann, als man will.
\begin{center}5.\end{center}
\noindent\hspace*{\parindent}\textls[120]{Erklärung}. Wenn in einer Reihe der vierten Art die letzte obere Grenze und die letzte untere einander gleich sind, so nennt man sie die \textls[120]{absolute} Grenze der Reihe.
\begin{center}6.\end{center}
\noindent\hspace*{\parindent}\textls[120]{Lehrsatz}. Wenn \(a,\,a',\,a''\) etc. eine Reihe ist, die die absolute Grenze \(A\) hat; \(b,\,b',\,b'',\,b'''\) etc. eine Reihe, deren absolute Grenze \(B\), so hat die Reihe \(a+b,\,a'+b',\,a''+b''\) etc. die absolute Grenze \(A+B\).
\par\medskip\noindent\rule{.24\linewidth}{.4pt}\par\smallskip
{\small\noindent\hspace*{\parindent}\textsuperscript{*}) Die erstere Reihe wird auch eine obere Grenze \(L'\), die zweite eine untere \(M'\) haben, auf die aber nicht weiter Rücksicht genommen zu werden braucht.\par}
